
\documentclass{article}
\usepackage{colortbl}
\usepackage{makecell}
\usepackage{multirow}
\usepackage{supertabular}

\begin{document}

\newcounter{utterance}

\centering \large Interaction Transcript for game `cladder', experiment `full\_v1.5\_default', episode 5585 with Qwen3.6{-}27B.
\vspace{24pt}

{ \footnotesize  \setcounter{utterance}{1}
\setlength{\tabcolsep}{0pt}
\begin{supertabular}{c@{$\;$}|p{.15\linewidth}@{}p{.15\linewidth}p{.15\linewidth}p{.15\linewidth}p{.15\linewidth}p{.15\linewidth}}
    \# & \multicolumn{2}{c}{Player} && \multicolumn{2}{c}{Game Master} \\
    \hline

    \theutterance \stepcounter{utterance}  
    & & & \multicolumn{4}{p{0.6\linewidth}}{
        \cellcolor[rgb]{0.9,0.9,0.9}{
            \makecell[{{p{\linewidth}}}]{
                \texttt{\tiny{[P1$\langle$GM]}}
                \texttt{You are an expert in causal inference. The following question is not a typical commonsense query, but rather a meticulously designed question created by a professor specializing in causal inference, intended to assess the students' mastery of the course content.} \\
\texttt{Imagine a self{-}contained, hypothetical world with only the following conditions, and without any unmentioned factors or causal relationships: Gender has a direct effect on treatment and peanut allergy. Treatment has a direct effect on peanut allergy. For patients who are not male and not receiving treatment, the probability of being allergic to peanuts is 84\%. For patients who are not male and receiving treatment, the probability of being allergic to peanuts is 71\%. For patients who are male and not receiving treatment, the probability of being allergic to peanuts is 21\%. For patients who are male and receiving treatment, the probability of being allergic to peanuts is 3\%. The overall probability of male gender is 54\%. For patients receiving treatment, would it be more likely to see being allergic to peanuts if the patient had received no treatment?} \\
\texttt{Start your answer with 'Yes' or 'No', followed by additional reasoning or evidence to support your explanation.} \\
            }
        }
    }
     \\ \\

    \theutterance \stepcounter{utterance}  
    & \multicolumn{4}{p{0.6\linewidth}}{
        \cellcolor[rgb]{0.561,0.737,0.561}{
            \makecell[{{p{\linewidth}}}]{
                \texttt{\tiny{[P1$\rangle$GM]}}
                \texttt{No} \\
\\ 
\texttt{**Reasoning:**} \\
\\ 
\texttt{To answer this question, we must determine the causal effect of treatment on peanut allergy for the specific subpopulation of patients who are receiving treatment. This is known as the **Conditional Average Treatment Effect (CATE)** or, more specifically, the effect of treatment on the treated (ETT) if we were comparing their actual outcome to their counterfactual outcome under no treatment. However, the question asks: "For patients receiving treatment, would it be more likely to see being allergic to peanuts if the patient had received no treatment?"} \\
\\ 
\texttt{This is asking for the comparison between:} \\
\texttt{1.  The observed probability of allergy for those receiving treatment: \$P(Y=1 | T=1)\$} \\
\texttt{2.  The counterfactual probability of allergy for those receiving treatment, had they not received treatment: \$P(Y\_\{T=0\}=1 | T=1)\$} \\
\\ 
\texttt{If \$P(Y\_\{T=0\}=1 | T=1) > P(Y=1 | T=1)\$, then the answer is "Yes" (it would be more likely to be allergic if they had not received treatment). If \$P(Y\_\{T=0\}=1 | T=1) < P(Y=1 | T=1)\$, the answer is "No".} \\
\\ 
\texttt{**Step 1: Identify the Causal Structure**} \\
\texttt{The problem states:} \\
\texttt{*   Gender (\$G\$) affects Treatment (\$T\$).} \\
\texttt{*   Gender (\$G} \\
            }
        }
    }
    & & \\ \\

    \theutterance \stepcounter{utterance}  
    & & & \multicolumn{2}{p{0.3\linewidth}}{
        \cellcolor[rgb]{0.9,0.9,0.9}{
            \makecell[{{p{\linewidth}}}]{
                \texttt{\tiny{[GM$|$GM]}}
                \texttt{No} \\
            }
        }
    }
    & & \\ \\

    \theutterance \stepcounter{utterance}  
    & & & \multicolumn{2}{p{0.3\linewidth}}{
        \cellcolor[rgb]{0.9,0.9,0.9}{
            \makecell[{{p{\linewidth}}}]{
                \texttt{\tiny{[GM$|$GM]}}
                \texttt{yes} \\
            }
        }
    }
    & & \\ \\

    \theutterance \stepcounter{utterance}  
    & & & \multicolumn{2}{p{0.3\linewidth}}{
        \cellcolor[rgb]{0.9,0.9,0.9}{
            \makecell[{{p{\linewidth}}}]{
                \texttt{\tiny{[GM$|$GM]}}
                \texttt{game\_result = LOSE} \\
            }
        }
    }
    & & \\ \\

\end{supertabular}
}

\end{document}
